
\documentclass{article}
\usepackage{colortbl}
\usepackage{makecell}
\usepackage{multirow}
\usepackage{supertabular}

\begin{document}

\newcounter{utterance}

\centering \large Interaction Transcript for game `cladder', experiment `full\_v1.5\_default', episode 4422 with reflexion{-}Qwen3.5{-}9B.
\vspace{24pt}

{ \footnotesize  \setcounter{utterance}{1}
\setlength{\tabcolsep}{0pt}
\begin{supertabular}{c@{$\;$}|p{.15\linewidth}@{}p{.15\linewidth}p{.15\linewidth}p{.15\linewidth}p{.15\linewidth}p{.15\linewidth}}
    \# & \multicolumn{2}{c}{Player} && \multicolumn{2}{c}{Game Master} \\
    \hline

    \theutterance \stepcounter{utterance}  
    & & & \multicolumn{4}{p{0.6\linewidth}}{
        \cellcolor[rgb]{0.9,0.9,0.9}{
            \makecell[{{p{\linewidth}}}]{
                \texttt{\tiny{[P1$\langle$GM]}}
                \texttt{You are an expert in causal inference. The following question is not a typical commonsense query, but rather a meticulously designed question created by a professor specializing in causal inference, intended to assess the students' mastery of the course content.} \\
\texttt{Imagine a self{-}contained, hypothetical world with only the following conditions, and without any unmentioned factors or causal relationships: Gender has a direct effect on smoking and peanut allergy. Smoking has a direct effect on tar deposit. Tar deposit has a direct effect on peanut allergy. Gender is unobserved. For nonsmokers, the probability of high tar deposit is 18\%. For smokers, the probability of high tar deposit is 64\%. For nonsmokers and with no tar deposit, the probability of being allergic to peanuts is 46\%. For nonsmokers and with high tar deposit, the probability of being allergic to peanuts is 77\%. For smokers and with no tar deposit, the probability of being allergic to peanuts is 17\%. For smokers and with high tar deposit, the probability of being allergic to peanuts is 43\%. The overall probability of smoking is 49\%. Will smoking decrease the chance of being allergic to peanuts?} \\
\texttt{Start your answer with 'Yes' or 'No', followed by additional reasoning or evidence to support your explanation.} \\
            }
        }
    }
     \\ \\

    \theutterance \stepcounter{utterance}  
    & \multicolumn{4}{p{0.6\linewidth}}{
        \cellcolor[rgb]{0.561,0.737,0.561}{
            \makecell[{{p{\linewidth}}}]{
                \texttt{\tiny{[P1$\rangle$GM]}}
                \texttt{No} \\
\\ 
\texttt{**Reasoning:**} \\
\\ 
\texttt{To determine if smoking decreases the chance of being allergic to peanuts, we must calculate the causal effect of smoking on peanut allergy. In this hypothetical world, the causal structure is defined as:} \\
\texttt{\$Gender \textbackslash{}rightarrow Smoking\$} \\
\texttt{\$Gender \textbackslash{}rightarrow Peanut Allergy\$} \\
\texttt{\$Smoking \textbackslash{}rightarrow Tar Deposit\$} \\
\texttt{\$Tar Deposit \textbackslash{}rightarrow Peanut Allergy\$} \\
\\ 
\texttt{Since **Gender** is an unobserved confounder that affects both the treatment (Smoking) and the outcome (Peanut Allergy), the observed association between smoking and peanut allergy will be biased. To answer the causal question "Will smoking decrease the chance...", we must estimate the **Average Causal Effect (ACE)**, which requires adjusting for the confounder (Gender) or using a method like backdoor adjustment. However, since the specific probabilities for Gender are not provided, we must rely on the **front{-}door criterion** or calculate the effect by conditioning on the mediator (Tar Deposit) to block the backdoor path through Gender.} \\
\\ 
\texttt{The front{-}door criterion states that if a variable \$M\$ (Tar Deposit) satisfies three conditions, we can identify the causal effect of \$X\$ (Smoking) on \$Y\$ (Peanut Allergy):} \\
\texttt{1.  \$X\$ affects \$M\$ (Smoking \$\textbackslash{}rightarrow\$ Tar Deposit).} \\
\texttt{2.  \$M\$ affects \$Y\$ (Tar Deposit \$\textbackslash{}rightarrow\$ Peanut Allergy).} \\
\texttt{3.  All paths from \$X\$ to \$Y\$ not going through \$M\$ are blocked by a set of observed variables. Here, the backdoor path is \$Smoking \textbackslash{}leftarrow Gender \textbackslash{}rightarrow Peanut Allergy\$. Since we condition on \$M\$ (Tar Deposit), does this block the path? No, conditioning on \$M\$ blocks the direct path from Smoking to Allergy via Tar, but the backdoor path \$Smoking \textbackslash{}leftarrow Gender \textbackslash{}rightarrow Allergy\$ remains open unless we condition on Gender.} \\
\\ 
\texttt{However, the standard approach when a confounder is unobserved but a valid mediator exists is to use the **Front{-}Door Formula**:} \\
\texttt{\$P(Y | do(X)) = \textbackslash{}sum\_m P(Y | M=m, X) P(M=m | X)\$} \\
\texttt{*Wait, this formula is incorrect for the front{-}door adjustment.* The correct front{-}door adjustment formula is:} \\
\texttt{\$P(Y | do(X)) = \textbackslash{}sum\_m \textbackslash{}sum\_\{x'\} P(Y | M=m, X=x') P(M=m | X=x') P(X=x')\$} \\
\texttt{Actually, the standard front{-}door formula is:} \\
\texttt{\$P(Y | do(X=x)) = \textbackslash{}sum\_m P(Y | M=m, X=x) \textbackslash{}sum\_\{x'\} P(M=m | X=x') P(X=x')\$} \\
\texttt{Let's re{-}verify the front{-}door adjustment logic.} \\
\texttt{The effect of \$X\$ on \$Y\$ is identified if we can sum over} \\
            }
        }
    }
    & & \\ \\

    \theutterance \stepcounter{utterance}  
    & & & \multicolumn{2}{p{0.3\linewidth}}{
        \cellcolor[rgb]{0.9,0.9,0.9}{
            \makecell[{{p{\linewidth}}}]{
                \texttt{\tiny{[GM$|$GM]}}
                \texttt{No} \\
            }
        }
    }
    & & \\ \\

    \theutterance \stepcounter{utterance}  
    & & & \multicolumn{2}{p{0.3\linewidth}}{
        \cellcolor[rgb]{0.9,0.9,0.9}{
            \makecell[{{p{\linewidth}}}]{
                \texttt{\tiny{[GM$|$GM]}}
                \texttt{no} \\
            }
        }
    }
    & & \\ \\

    \theutterance \stepcounter{utterance}  
    & & & \multicolumn{2}{p{0.3\linewidth}}{
        \cellcolor[rgb]{0.9,0.9,0.9}{
            \makecell[{{p{\linewidth}}}]{
                \texttt{\tiny{[GM$|$GM]}}
                \texttt{game\_result = WIN} \\
            }
        }
    }
    & & \\ \\

\end{supertabular}
}

\end{document}
